% Full stage induction and finite Boolean normalization for the concluding Kamp construction.

\section{Kamp's interval components and their Boolean normalization}
\label{sec:annex-p09-kamp-iii2-lemma10}

\begin{annexitem}{annex:p9.kamp.iii2.lemma10}{thm:p09-kamp-iii2-lemma10}{Both Bet constructors and every endpoint conjunction}
\begin{proof}
Induct on the supplied finite stage of the exact source Bet recursion,
simultaneously for all pairs of distinct endpoint names. Stage zero has
no members, so its assertion is vacuous. At stage one, an existential
matrix contains only At conditions at its quantified point. Their
conjunction is an elementary label, including the source truth label
for an empty list. The word
\((T^K,T^K,\theta,T^K,T^K)\) represents its strict interior
witness condition. A universal stage-one matrix is a disjunction of At
conditions, an elementary label \(\eta\), including falsehood for
an empty list. Its representative is \((T^K,\eta,T^K)\).
Thus both bases have actual elementary representations even when their
open intervals are empty.

Suppose all components of stage \(n\) have their representations.
A retained component needs no new construction. For a newly constructed
existential component, every matrix conjunct belongs to
\(\operatorname{At}(z)\), \(\operatorname{Bet}_n(x,z)\),
or \(\operatorname{Bet}_n(z,y)\). Group it accordingly.
The first group has an elementary conjunction at \(z\);
the other two groups have supplied interval properties by induction.
\Cref{thm:p09-kamp-iii2-lemma5} applies with exactly the original
strict guard \(x<z<y\) and produces the new interval representation.
Every conjunct and every pair of elementary disjuncts is retained in
that construction. Its empty-family convention retains an actual
strict midpoint witness.

For a newly constructed universal component the three corresponding
groups occur disjunctively. Apply
\cref{thm:p09-kamp-iii2-lemma9}: the At disjunction has an elementary
history and the two Bet groups have their inductively supplied
representations. The universal guard is unchanged. Lemma~9's explicit
empty cases ensure that an empty matrix still expresses absence of
interior points. This proves the induction step for both constructors.

Every Bet component belongs to some finite stage, so it has property
\(I\). Any finite list of such components has finitely many
given representations. Each At formula at an outer endpoint belongs
to \(\operatorname{E}\) at that endpoint, and the endpoint
insertion theorem gives its \(I(x,y)\) presentation. Repeatedly
apply the conjunction fusion theorem to the finite list, distributing
over all pairs at each step. This gives the claimed finite conjunction
representation. For no factors use \((T^K,T^K,T^K)\), true on
every ordered pair. Each stage operation is finite and uses only
earlier-stage constructions, proving termination.

Every underlying elementary formula already has its own interval
property by the definition of \(I\); no additional expressibility
theorem is needed to treat it as one disjunct. The only complete-order
dependencies in this induction are the actual interval representations
and Lemmas~5 and 9; their scope is never promoted to \(\mathsf{Lin}\).

For the example, a two-interval truth word has an internal witness
exactly when the open interval is nonempty. The one-interval word with
false interval label holds exactly when that open interval is empty.
Their endpoint truth tests add no restriction. This proves the
complement relationship and the adjacent-endpoint case directly
\cite[ch.~III, lemma~10, p.~93]{Kamp1968}.
\end{proof}
\end{annexitem}

\begin{annexitem}{annex:p9.kamp.interval.boolean.normalization}{thm:p09-kamp-interval-boolean-normalization}{Every guarded negation and both finite normal-form conversions}
\begin{proof}
We give a finite construction for the complement of an At or Bet
component. General Boolean expressions are then put in negation normal
form and each complemented component is replaced by that construction.
The resulting external expression uses only conjunction and disjunction.

An At component is a finite conjunction of unary literals or the
particular source tautology. Negating a literal changes its sign;
negating a conjunction gives the disjunction of its complemented
literals, each an At component. Negating the particular tautology
gives the source falsehood. These are valid identities for every
unary interpretation and do not introduce order atoms. The truth
and falsehood conventions also cover empty Boolean lists.

Now proceed by induction on a finite construction history of a Bet
component at endpoints \((x,y)\). A retention step passes to the
earlier history. Write \(g(z)\) for \(x<z<y\).
If the component is
\(\exists z\,(g(z)\land\bigwedge_i\theta_i)\), its negation is
\[
                  \forall z\,(g(z)\to\bigvee_i\neg\theta_i).
 \tag{1}
\]
Each \(\theta_i\) is At at \(z\), or a proper earlier
Bet component on \((x,z)\) or \((z,y)\). Replace its negation
by the already constructed positive expression for that component.
The matrix of (1) is now a positive Boolean combination of leaves of
the following five kinds:
\[
 \operatorname{At}(x),\quad\operatorname{At}(y),\quad
 \operatorname{At}(z),\quad\operatorname{Bet}(x,z),\quad
 \operatorname{Bet}(z,y).
\]
There are no additional free variables. Put this finite matrix in
conjunctive normal form by distributivity; the leaves are treated as
propositional inputs for this syntactic operation. Distribute the
universal quantifier over the resulting finite conjunction of clauses.

In one clause, separate the disjunction \(A(x,y)\) of endpoint
At leaves from the disjunction \(B(z)\) of the other three
kinds. The first has a fixed truth value while \(z\) varies.
Therefore
\[
 \forall z\,(g(z)\to A(x,y)\lor B(z))
 \longleftrightarrow A(x,y)\lor\forall z\,(g(z)\to B(z)).
 \tag{2}
\]
If \(A\) is true both sides are true. If it is false,
both sides are the same guarded universal. This also proves (2)
when the guarded interval is empty. When no local leaves occur,
use \(B(z)=F^K_z\); the universal component then asserts
that the interval is empty and cannot be deleted. The last
universal in (2) is a genuine Bet component: take the maximum
of the finitely many stages of its local Bet leaves and then
the next source universal stage. An empty set of such leaves
uses the stage-one falsehood case. Combining the clauses gives
a positive Boolean expression over the required outer-endpoint
classes.

If instead the original component is
\(\forall z\,(g(z)\to\bigvee_i\theta_i)\), its negation is
\[
                  \exists z\,(g(z)\land\bigwedge_i\neg\theta_i).
 \tag{3}
\]
Replace each complemented constituent by its inductively constructed
positive expression. Put the resulting finite matrix in disjunctive
normal form and distribute the existential quantifier over all its
terms. In a term, separate the endpoint conjunction \(A(x,y)\)
from the conjunction \(B(z)\) of local At and Bet leaves.
Now
\[
 \exists z\,(g(z)\land A(x,y)\land B(z))
 \longleftrightarrow A(x,y)\land\exists z\,(g(z)\land B(z)).
 \tag{4}
\]
If \(A\) is false both sides are false. If it is true,
both sides require the same actual witness. For an empty local
conjunction use \(T^K_z\); the remaining existential still
requires a point strictly between \(x\) and \(y\).
The existential in (4) is a Bet component at one stage above
the maximum stage of its local leaves, or the stage-one truth
case when no such leaf occurs. Combining all terms gives the
required positive expression.

The finite Boolean conversions preserve false empty disjunctions and
true empty conjunctions through the declared source constants. All
bound variables are made fresh before the substitutions and distributions.
The algorithm recurses only on proper original constituents. It does
not restart on the new Bet components produced by (2) or (4).
Consequently each recursive branch terminates at an original At
leaf or a lower construction stage. All normal-form conversions
are finite. These arguments prove equivalence over every linear
structure, indeed by the displayed Boolean and quantifier identities,
without assuming order completeness.

Finally every resulting Bet leaf has property \(I(x,y)\)
by Lemma~10, and each outer At leaf has it by endpoint insertion.
Finite conjunction uses fusion and finite disjunction concatenates
the supplied elementary disjunct lists. Thus the whole positive
expression has an actual interval representation over
\(\mathsf{Com}\). This is a justified closure for the
specified source interval expressions; it does not assume closure
under negation for arbitrary formulas already bearing property \(I\).
It supplies the detailed finite transformations underlying source
p.~62 \cite[ch.~III, pp.~61--62, 93]{Kamp1968}.
\end{proof}
\end{annexitem}
